\documentclass[aspectratio=169,11pt]{beamer}

% Packages
\usepackage[utf8]{inputenc}
\usepackage[T1]{fontenc}
\usepackage{lmodern}
\usepackage{microtype}
\usepackage{amsmath,amssymb}
\usepackage{booktabs}
\usepackage{graphicx}
\usepackage{tikz}
\usepackage{pgfplots}
\usepackage{xcolor}
\usepackage{ragged2e}
\usepackage{colortbl}
\usepackage{array}
\usepackage{hyperref}
\usepackage{adjustbox}

\usetikzlibrary{shapes.geometric, arrows.meta, positioning, calc, backgrounds, decorations.pathreplacing, shadows, fadings}
\pgfplotsset{compat=1.18}

% --- Color Palette: Warm Professional ---
\definecolor{DeepNavy}{HTML}{2E4057}
\definecolor{Teal}{HTML}{048A81}
\definecolor{WarmOrange}{HTML}{E85D04}
\definecolor{SoftPurple}{HTML}{9D4EDD}
\definecolor{WarmGray}{HTML}{6C757D}
\definecolor{LightGray}{HTML}{E9ECEF}
\definecolor{Cream}{HTML}{FBF8F1}

\setbeamercolor{background canvas}{bg=Cream}
\setbeamercolor{normal text}{fg=DeepNavy}
\setbeamercolor{frametitle}{fg=Cream, bg=DeepNavy}
\setbeamercolor{title}{fg=Cream, bg=DeepNavy}
\setbeamercolor{item}{fg=Teal}

\begin{document}

\begin{frame}{Granular Instrumental Variables isolate idiosyncratic shocks}
    \vspace{0.2cm}
    \centering
    The GIV $z_t$ is constructed from observables, $y_{it}$, using weights $\Gamma \in \mathbb{R}^N$ orthogonal to factor loadings:
    
    \vspace{0.1cm}
    \begin{equation*}
        z_t := \Gamma' y_t = \sum_{i=1}^{N} \Gamma_i y_{it}
    \end{equation*}
    
    \vspace{0.1cm}
    By construction, this isolates a linear combination of idiosyncratic shocks $u_t$:
    
    \vspace{0.1cm}
    \begin{equation*}
        z_t = \Gamma' u_t \quad \Rightarrow \quad \text{exogenous to aggregate shocks } \eta_t, \varepsilon_t
    \end{equation*}

    \vspace{0.2cm}
    This allows us to construct the sample estimators:
    \vspace{0.1cm}
    
    \adjustbox{max width=\textwidth}{%
    $
    \displaystyle
    \hat{\psi}_T = \frac{\sum_{t=1}^T p_t z_t}{\sum_{t=1}^T y_{St} z_t} 
    \hspace{2cm} 
    \hat{\phi}^d_T = \frac{\sum_{t=1}^T y_{\tilde{E}t} z_t}{\sum_{t=1}^T p_t z_t}
    $
    }
\end{frame}

\begin{frame}{Moment conditions link structural equations to the GIV}
    \vspace{0.5cm}
    \centering
    The estimators follow directly from the exogeneity of $z_t$ to structural shocks:
    
    \vspace{0.6cm}
    \begin{minipage}{0.48\textwidth}
        \centering
        \textbf{Inverse Supply Elasticity}
        
        \vspace{0.2cm}
        From $p_t = \psi y_{St} + \dots + \varepsilon_t$:
        \begin{equation*}
            \mathbb{E}[(p_t - \psi y_{St}) z_t] = 0
        \end{equation*}
        \vspace{0.1cm}
        \begin{equation*}
            \Rightarrow \quad \psi = \frac{\mathbb{E}[p_t z_t]}{\mathbb{E}[y_{St} z_t]}
        \end{equation*}
    \end{minipage}
    \hfill
    \begin{minipage}{0.48\textwidth}
        \centering
        \textbf{Demand Elasticity}
        
        \vspace{0.2cm}
        For weights $\tilde{E}$ where $\mathbb{E}[u_{\tilde{E}t} u_{\Gamma t}] = 0$:
        \begin{equation*}
            \mathbb{E}[(y_{\tilde{E}t} - \phi^d p_t) z_t] = 0
        \end{equation*}
        \vspace{0.1cm}
        \begin{equation*}
            \Rightarrow \quad \phi^d = \frac{\mathbb{E}[y_{\tilde{E}t} z_t]}{\mathbb{E}[p_t z_t]}
        \end{equation*}
    \end{minipage}
\end{frame}

\begin{frame}{Proposition 1: The GIV estimator is asymptotically consistent}
    \vspace{0.5cm}
    \centering
    Under standard regularity conditions, as the number of time periods $T \to \infty$:
    
    \vspace{0.8cm}
    \begin{minipage}{0.45\textwidth}
        \centering
        \textbf{Inverse Supply Elasticity}
        
        \vspace{0.4cm}
        $ \hat{\psi}_T \xrightarrow{\text{a.s.}} \psi $
    \end{minipage}
    \hfill
    \begin{minipage}{0.45\textwidth}
        \centering
        \textbf{Demand Elasticity}
        
        \vspace{0.4cm}
        $ \hat{\phi}^d_T \xrightarrow{\text{a.s.}} \phi^d $ \quad (if $\psi \neq 0$)
    \end{minipage}

    \vspace{1cm}
    \textbf{Intuition:} The orthogonal weighting $\Gamma$ ensures the instrument is uncorrelated with the structural error term, resolving simultaneity bias.
\end{frame}

\begin{frame}{Pass-through multipliers capture shock amplification}
    \vspace{0.5cm}
    \centering
    To understand the variance of these estimators, we define two key multipliers:
    
    \vspace{0.8cm}
    \begin{minipage}{0.45\textwidth}
        \centering
        \textbf{Quantity Pass-through}
        
        \vspace{0.3cm}
        \begin{equation*}
            M := \frac{1}{1 - \psi\phi^d}
        \end{equation*}
    \end{minipage}
    \hfill
    \begin{minipage}{0.45\textwidth}
        \centering
        \textbf{Price Pass-through}
        
        \vspace{0.3cm}
        \begin{equation*}
            \mu := \psi M
        \end{equation*}
    \end{minipage}

    \vspace{1cm}
    These represent the pass-through from idiosyncratic and aggregate shocks to the aggregate quantity and the price, respectively.
\end{frame}

\begin{frame}{Proposition 2: Asymptotic variance scales with aggregate noise}
    \vspace{0.5cm}
    \centering
    The estimators converge in distribution to a normal distribution:
    
    \vspace{0.4cm}
    \adjustbox{max width=\textwidth}{%
    \begin{minipage}{0.8\textwidth}
        \begin{equation*}
            \sqrt{T}(\hat{\psi}_T - \psi) \xrightarrow{d} \mathcal{N}(0, \sigma^2_\psi), \quad \sigma_\psi = \frac{\mathbb{E}[u_{\Gamma t}^2]^{1/2}}{|\mathbb{E}[u_{St} u_{\Gamma t}]|} \frac{\sigma_{\varepsilon^\psi}}{|M|}
        \end{equation*}
    \end{minipage}
    }
    
    \vspace{0.6cm}
    \adjustbox{max width=\textwidth}{%
    \begin{minipage}{0.8\textwidth}
        \begin{equation*}
            \sqrt{T}(\hat{\phi}^d_T - \phi^d) \xrightarrow{d} \mathcal{N}(0, \sigma^2_{\phi^d}), \quad \sigma_{\phi^d} = \frac{\mathbb{E}[u_{\Gamma t}^2]^{1/2}}{|\mathbb{E}[u_{St} u_{\Gamma t}]|} \frac{\sigma_{\varepsilon^\phi}}{|\mu|}
        \end{equation*}
    \end{minipage}
    }
\end{frame}

\begin{frame}{High idiosyncratic volatility improves estimator precision}
    \vspace{0.5cm}
    \centering
    The variance equations reveal the signal-to-noise ratio inherent to GIV:
    
    \vspace{0.8cm}
    \begin{minipage}{0.3\textwidth}
        \centering
        \textcolor{Teal}{\textbf{Signal Strength}}
        
        \vspace{0.2cm}
        $|\mathbb{E}[u_{St} u_{\Gamma t}]|$
        
        \vspace{0.2cm}
        Larger granular shocks increase relevance.
    \end{minipage}
    \hfill
    \begin{minipage}{0.3\textwidth}
        \centering
        \textcolor{WarmOrange}{\textbf{Instrument Noise}}
        
        \vspace{0.2cm}
        $\mathbb{E}[u_{\Gamma t}^2]^{1/2}$
        
        \vspace{0.2cm}
        Variance of the constructed GIV itself.
    \end{minipage}
    \hfill
    \begin{minipage}{0.3\textwidth}
        \centering
        \textcolor{SoftPurple}{\textbf{Structural Noise}}
        
        \vspace{0.2cm}
        $\sigma_{\varepsilon} / |M|$
        
        \vspace{0.2cm}
        Unobserved aggregate shocks dampening precision.
    \end{minipage}
\end{frame}

\begin{frame}{Proposition 3: Optimal weights use an orthogonal projection}
    \vspace{0.5cm}
    \centering
    To minimize the asymptotic variances $\sigma^2_\psi$ and $\sigma^2_{\phi^d}$, we use the $N \times N$ projection matrix $Q$ that is orthogonal to the factor loadings $\lambda$:
    
    \vspace{0.3cm}
    \begin{equation*}
        Q := I - \lambda(\lambda'\lambda)^{-1}\lambda'
    \end{equation*}

    \vspace{0.5cm}
    The optimal weights vector $\Gamma^*$ is the projection of the size vector $S$:
    
    \vspace{0.3cm}
    \begin{equation*}
        \Gamma^* = S'Q
    \end{equation*}

    \vspace{0.5cm}
    \textbf{Result:} $\Gamma^*$ is the vector closest to $S$, while remaining orthogonal to the unobserved aggregate factor loadings.
\end{frame}

\end{document}